% !TeX root = main.tex

\setlength{\abovedisplayskip}{7pt}
\setlength{\belowdisplayskip}{7pt}
\setlength{\abovedisplayshortskip}{4pt}
\setlength{\belowdisplayshortskip}{4pt}

\tableofcontents
\newpage

\section{Цель работы}

Синтезировать дискретные стабилизирующий и следящий регуляторы
методом модального управления. Построить устройство оценки
пониженной размерности и получить законы управления при неполной
информации о состоянии объекта.

\section{Модель объекта управления}

\subsection{Исходные данные}

Для варианта 3 задан объект третьего типа с параметрами
\begin{equation}
  k_1=3{,}26,\qquad k_2=1,\qquad
  a_0^1=a_0^2=1,\qquad T_1=9{,}5,\qquad
  \xi=0{,}36,\qquad T=1\,\text{с}
  \label{eq:plant-data}
\end{equation}
Передаточная функция объекта имеет вид
\begin{equation}
  W_{\text{н}}(s)
  =\frac{k_1k_2}{T_1^2s^2+2T_1\xi s+a_0^1a_0^2}
  =\frac{3{,}26}{90{,}25s^2+6{,}84s+1}
  \label{eq:continuous-transfer}
\end{equation}
Соответствующее дифференциальное уравнение:
\begin{equation}
  90{,}25\ddot y+6{,}84\dot y+y=3{,}26u
  \label{eq:continuous-equation}
\end{equation}
Выбираем вектор состояния
\begin{equation}
  \bm x=\begin{pmatrix}y\\\dot y\end{pmatrix},
  \qquad y=C\bm x,\qquad C=\begin{pmatrix}1&0\end{pmatrix}
  \label{eq:state-choice}
\end{equation}
Непрерывная модель объекта записывается в виде
\begin{equation}
  \dot{\bm x}=A_{\text{н}}\bm x+B_{\text{н}}u,
  \qquad
  A_{\text{н}}=
  \begin{pmatrix}
    0&1\\
    -\dfrac{1}{90{,}25}&-\dfrac{6{,}84}{90{,}25}
  \end{pmatrix},
  \qquad
  B_{\text{н}}=\begin{pmatrix}0\\\dfrac{3{,}26}{90{,}25}\end{pmatrix}
  \label{eq:continuous-state-space}
\end{equation}

\subsection{Точная дискретизация}

При удержании управления на каждом интервале дискретизации получаем
\begin{equation}
  \bm x(k+1)=A\bm x(k)+Bu(k),
  \qquad y(k)=C\bm x(k)
  \label{eq:discrete-plant}
\end{equation}
Матрицы дискретной модели определяются выражениями
\begin{equation}
  A=e^{A_{\text{н}}T},
  \qquad
  B=\int_0^T e^{A_{\text{н}}\tau}B_{\text{н}}\,d\tau
    =A_{\text{н}}^{-1}(A-I_2)B_{\text{н}}
  \label{eq:discretization}
\end{equation}
Для заданного периода $T=1\,\text{с}$ получаем
\begin{equation}
  A\approx
  \begin{pmatrix}
    0{,}99460215&0{,}96126742\\
    -0{,}01065116&0{,}92174819
  \end{pmatrix},
  \qquad
  B\approx\begin{pmatrix}0{,}01759701\\0{,}03472279\end{pmatrix}
  \label{eq:discrete-matrices}
\end{equation}

\subsection{Проверка структурных свойств}

Матрица управляемости равна
\begin{equation}
  \mathcal C=\begin{pmatrix}B&AB\end{pmatrix}
  \approx
  \begin{pmatrix}
    0{,}01759701&0{,}05087991\\
    0{,}03472279&0{,}03181824
  \end{pmatrix}
  \label{eq:controllability}
\end{equation}
Её определитель отличен от нуля:
\begin{equation}
  \det\mathcal C\approx-0{,}00120679\ne0,
  \qquad \operatorname{rank}\mathcal C=2
  \label{eq:controllability-rank}
\end{equation}
Следовательно, пара $(A,B)$ полностью управляема.

Матрица наблюдаемости имеет вид
\begin{equation}
  \mathcal O=\begin{pmatrix}C\\CA\end{pmatrix}
  \approx
  \begin{pmatrix}
    1&0\\
    0{,}99460215&0{,}96126742
  \end{pmatrix}
  \label{eq:observability}
\end{equation}
Поскольку
\begin{equation}
  \det\mathcal O\approx0{,}96126742\ne0,
  \qquad \operatorname{rank}\mathcal O=2
  \label{eq:observability-rank}
\end{equation}
пара $(C,A)$ полностью наблюдаема.

\section{Синтез стабилизирующего регулятора}

\subsection{Эталонная модель}

Управление выбирается в виде линейной обратной связи по состоянию:
\begin{equation}
  u(k)=-K_{\text{с}}\bm x(k),
  \qquad F_{\text{с}}=A-BK_{\text{с}}
  \label{eq:stabilizing-control}
\end{equation}
Для оптимальной по быстродействию дискретной системы назначаем
оба желаемых корня в начале координат:
\begin{equation}
  z_1^*=z_2^*=0,
  \qquad D_{\text{с}}(z)=z^2
  \label{eq:stabilizing-polynomial}
\end{equation}
Матрицы эталонной модели выбираем следующим образом:
\begin{equation}
  \Gamma_{\text{с}}=\begin{pmatrix}0&1\\0&0\end{pmatrix},
  \qquad H_{\text{с}}=\begin{pmatrix}1&0\end{pmatrix}
  \label{eq:stabilizing-reference}
\end{equation}
Матрица наблюдаемости пары $(H_{\text{с}},\Gamma_{\text{с}})$ равна
\begin{equation}
  \begin{pmatrix}H_{\text{с}}\\H_{\text{с}}\Gamma_{\text{с}}\end{pmatrix}
  =I_2
  \label{eq:stabilizing-reference-observability}
\end{equation}

\subsection{Расчёт обратной связи}

В методе модального управления матрица обратной связи определяется
из уравнения Сильвестра:
\begin{equation}
  A M_{\text{с}}-M_{\text{с}}\Gamma_{\text{с}}=BH_{\text{с}},
  \qquad K_{\text{с}}=H_{\text{с}}M_{\text{с}}^{-1}
  \label{eq:stabilizing-sylvester}
\end{equation}
Представляя $M_{\text{с}}$ через столбцы, получаем
\begin{equation}
  A\bm m_1=B,
  \qquad A\bm m_2=\bm m_1,
  \qquad
  M_{\text{с}}=\begin{pmatrix}A^{-1}B&A^{-2}B\end{pmatrix}
  \label{eq:stabilizing-sylvester-columns}
\end{equation}
Численное решение имеет вид
\begin{equation}
  M_{\text{с}}\approx
  \begin{pmatrix}
    -0{,}01850881&-0{,}05724458\\
    0{,}03745670&0{,}03997511
  \end{pmatrix},
  \qquad \det M_{\text{с}}\approx0{,}00140430\ne0
  \label{eq:stabilizing-sylvester-numeric}
\end{equation}
Следовательно,
\begin{equation}
  K_{\text{с}}\approx\begin{pmatrix}28{,}46618609&40{,}76373719\end{pmatrix}
  \label{eq:stabilizing-gain}
\end{equation}
Закон управления принимает вид
\begin{equation}
  u(k)\approx-28{,}46618609x_1(k)-40{,}76373719x_2(k)
  \label{eq:stabilizing-control-numeric}
\end{equation}
Матрица замкнутой системы равна
\begin{equation}
  F_{\text{с}}\approx
  \begin{pmatrix}
    0{,}49368248&0{,}24394766\\
    -0{,}99907656&-0{,}49368248
  \end{pmatrix}
  \label{eq:stabilizing-closed-loop}
\end{equation}

\section{Синтез следящего регулятора}

\subsection{Генератор задающего воздействия}

По данным варианта 3 параметры задания равны
\begin{equation}
  g_0=2{,}8,\qquad g_1=0,\qquad A_g=0,\qquad \omega_g=0
  \label{eq:reference-data}
\end{equation}
Таким образом, требуется воспроизводить постоянный сигнал
\begin{equation}
  g(k)=2{,}8
  \label{eq:reference-signal}
\end{equation}
Дискретный генератор этого сигнала имеет описание
\begin{equation}
  \zeta(k+1)=\zeta(k),
  \qquad g(k)=\zeta(k),
  \qquad \zeta(0)=2{,}8
  \label{eq:reference-generator}
\end{equation}

\subsection{Расширенная модель объекта}

Определим ошибку слежения и встроенную модель постоянного задания:
\begin{equation}
  e(k)=g(k)-y(k),
  \qquad \eta(k+1)=\eta(k)+e(k)
  \label{eq:internal-model}
\end{equation}
Для расширенного вектора состояния получаем
\begin{equation}
  \bar{\bm x}(k)=\begin{pmatrix}\eta(k)\\\bm x(k)\end{pmatrix},
  \qquad
  \bar{\bm x}(k+1)
    =\bar A\bar{\bm x}(k)+\bar B u(k)+\bar B_g g(k)
  \label{eq:augmented-plant}
\end{equation}
Матрицы расширенной модели имеют вид
\begin{equation}
  \bar A=\begin{pmatrix}1&-C\\0&A\end{pmatrix},
  \qquad
  \bar B=\begin{pmatrix}0\\B\end{pmatrix},
  \qquad
  \bar B_g=\begin{pmatrix}1\\0\\0\end{pmatrix}
  \label{eq:augmented-matrices}
\end{equation}
Проверим управляемость расширенной системы:
\begin{equation}
  \bar{\mathcal C}
    =\begin{pmatrix}\bar B&\bar A\bar B&\bar A^2\bar B\end{pmatrix},
  \qquad
  \det\bar{\mathcal C}\approx4{,}19417256\cdot10^{-5}\ne0
  \label{eq:augmented-controllability}
\end{equation}
Следовательно, её ранг равен трём и можно назначить все три корня
характеристического полинома.

\subsection{Расчёт коэффициентов регулятора}

Для расширенной системы выбираем
\begin{equation}
  D_{\text{сл}}(z)=z^3,
  \qquad
  \Gamma_{\text{сл}}=
  \begin{pmatrix}0&1&0\\0&0&1\\0&0&0\end{pmatrix},
  \qquad
  H_{\text{сл}}=\begin{pmatrix}1&0&0\end{pmatrix}
  \label{eq:tracking-reference}
\end{equation}
Для выбранной эталонной модели матрица наблюдаемости является
единичной матрицей третьего порядка. Матрицу обратной связи
находим из уравнений
\begin{equation}
  \bar A M_{\text{сл}}-M_{\text{сл}}\Gamma_{\text{сл}}
    =\bar B H_{\text{сл}},
  \qquad \bar K=H_{\text{сл}}M_{\text{сл}}^{-1}
  \label{eq:tracking-sylvester}
\end{equation}
Аналогично расчёту~\eqref{eq:stabilizing-sylvester-columns} получаем
\begin{equation}
  M_{\text{сл}}=
  \begin{pmatrix}
    \bar A^{-1}\bar B&\bar A^{-2}\bar B&\bar A^{-3}\bar B
  \end{pmatrix}
  \label{eq:tracking-sylvester-solution}
\end{equation}
Численное решение равно
\begin{equation}
  M_{\text{сл}}\approx
  \begin{pmatrix}
    -0{,}01850881&-0{,}07575339&-0{,}17412527\\
    -0{,}01850881&-0{,}05724458&-0{,}09837188\\
    0{,}03745670&0{,}03997511&0{,}04223207
  \end{pmatrix}
  \label{eq:tracking-sylvester-numeric}
\end{equation}
Определитель этой матрицы отличен от нуля:
\begin{equation}
  \det M_{\text{сл}}\approx-5{,}26491235\cdot10^{-5}\ne0
  \label{eq:tracking-sylvester-determinant}
\end{equation}
Матрица обратной связи равна
\begin{equation}
  \bar K=\begin{pmatrix}-K_\eta&\kappa_1&\kappa_2\end{pmatrix}
  \approx
  \begin{pmatrix}-28{,}77293455&71{,}44381439&47{,}78282739\end{pmatrix}
  \label{eq:tracking-gain}
\end{equation}
Следящий регулятор записывается в виде
\begin{equation}
  \begin{aligned}
    \eta(k+1)&=\eta(k)+g(k)-y(k),\\
    u(k)&=\kappa_1g(k)+K_\eta\eta(k)
      -\kappa_1x_1(k)-\kappa_2x_2(k)
  \end{aligned}
  \label{eq:tracking-controller}
\end{equation}
После подстановки численных коэффициентов получаем
\begin{equation}
  u(k)\approx71{,}44381439\bigl(g(k)-y(k)\bigr)
    +28{,}77293455\eta(k)-47{,}78282739x_2(k)
  \label{eq:tracking-control-numeric}
\end{equation}

\section{Регуляторы при неполной информации}

\subsection{Устройство оценки пониженной размерности}

Считаем доступным для измерения только выход объекта.
Для нечётного варианта 3 строим устройство оценки размерности
$n-1=1$, то есть скалярный наблюдатель.
Выбираем матрицы его описания
\begin{equation}
  F_{\text{н}}=0,\qquad G=1
  \label{eq:observer-dynamics}
\end{equation}
Матрицу преобразования находим из уравнения
\begin{equation}
  M_{\text{н}}A-F_{\text{н}}M_{\text{н}}=GC
  \label{eq:observer-sylvester}
\end{equation}
При выбранных коэффициентах получаем
\begin{equation}
  M_{\text{н}}=CA^{-1}
  \approx\begin{pmatrix}0{,}99432245&-1{,}03695324\end{pmatrix},
  \qquad
  M_{\text{н}}B\approx-0{,}01850881
  \label{eq:observer-matrices}
\end{equation}
Уравнение устройства оценки принимает вид
\begin{equation}
  \begin{aligned}
    \hat z(k+1)&=F_{\text{н}}\hat z(k)+Gy(k)+M_{\text{н}}Bu(k)\\
      &\approx y(k)-0{,}01850881u(k)
  \end{aligned},
  \qquad \hat z(0)=0
  \label{eq:observer}
\end{equation}
Для невязки $\tilde z(k)=\hat z(k)-M_{\text{н}}\bm x(k)$ имеем
\begin{equation}
  \tilde z(k+1)=F_{\text{н}}\tilde z(k)=0
  \label{eq:observer-error}
\end{equation}
Вектор состояния восстанавливается по измерению и состоянию наблюдателя:
\begin{equation}
  \hat{\bm x}(k)
  =\begin{pmatrix}C\\M_{\text{н}}\end{pmatrix}^{-1}
    \begin{pmatrix}y(k)\\\hat z(k)\end{pmatrix}
  \label{eq:state-estimate}
\end{equation}
Матрица восстановления невырождена:
\begin{equation}
  \det\begin{pmatrix}C\\M_{\text{н}}\end{pmatrix}
    \approx-1{,}03695324\ne0
  \label{eq:reconstruction-determinant}
\end{equation}
Таким образом,
\begin{equation}
  \hat x_1(k)=y(k),
  \qquad
  \hat x_2(k)\approx
    \frac{\hat z(k)-0{,}99432245y(k)}{-1{,}03695324}
  \label{eq:state-estimate-numeric}
\end{equation}

\subsection{Стабилизация с устройством оценки}

В законе~\eqref{eq:stabilizing-control} заменяем состояние его оценкой:
\begin{equation}
  u(k)=-K_{\text{с}}\hat{\bm x}(k)
      =-E_{\text{с}}y(k)-D_{\text{с}}\hat z(k)
  \label{eq:estimated-stabilizing-control}
\end{equation}
Коэффициенты регулятора определяются соотношением
\begin{equation}
  \begin{pmatrix}E_{\text{с}}&D_{\text{с}}\end{pmatrix}
    =K_{\text{с}}\begin{pmatrix}C\\M_{\text{н}}\end{pmatrix}^{-1}
    \approx\begin{pmatrix}67{,}55406166&-39{,}31106609\end{pmatrix}
  \label{eq:estimated-stabilizing-gains}
\end{equation}
Для реализации получаем
\begin{equation}
  \begin{aligned}
    \hat z(k+1)&\approx y(k)-0{,}01850881u(k),\\
    u(k)&\approx-67{,}55406166y(k)+39{,}31106609\hat z(k)
  \end{aligned}
  \label{eq:estimated-stabilizing-controller}
\end{equation}

\subsection{Слежение с устройством оценки}

В следящем регуляторе используется то же устройство оценки.
Состояние встроенной модели формируется внутри регулятора
по измеренной ошибке слежения. Коэффициенты при измерении
и состоянии наблюдателя равны
\begin{equation}
  \begin{pmatrix}E_{\text{сл}}&D_{\text{сл}}\end{pmatrix}
    =\begin{pmatrix}\kappa_1&\kappa_2\end{pmatrix}
      \begin{pmatrix}C\\M_{\text{н}}\end{pmatrix}^{-1}
    \approx\begin{pmatrix}117{,}26221428&-46{,}08002148\end{pmatrix}
  \label{eq:estimated-tracking-gains}
\end{equation}
Итоговые уравнения регулятора:
\begin{equation}
  \begin{aligned}
    \eta(k+1)&=\eta(k)+g(k)-y(k),\\
    \hat z(k+1)&\approx y(k)-0{,}01850881u(k),\\
    u(k)&\approx71{,}44381439g(k)+28{,}77293455\eta(k)\\
      &\hphantom{{}\approx{}}-117{,}26221428y(k)+46{,}08002148\hat z(k)
  \end{aligned}
  \label{eq:estimated-tracking-controller}
\end{equation}
